% 两层DSE架构
\section{两层DSE架构：分离"选型"与"判断"}

在进入具体约束之前，先建立框架的整体结构。我们将DSE流程分为两层。

\textbf{外层：技术选型。}
外层决定离散的架构选择——拓扑参数（Dragonfly的$a,p,h$）、互连标准（UCIe等级、SerDes速率）、
裸片参数（面积$A$、TDP功耗$P$、bump pitch）、冷却方案、裸片布局。
外层是一个枚举/搜索问题：在离散的、可能带上非线性经典约束的候选空间中产生设计点。
外层策略可以是暴力枚举、启发式剪枝或基于agent的智能选型——取决于设计阶段和计算预算。

\textbf{内层：多约束耦合判断。}
内层的任务是在外层固定的技术选型下，回答三个问题：
\begin{enumerate}
    \item 拓扑能否承载目标交换带宽？
    \item 物理资源（bump、布线容量、散热能力）是否足够？
    \item 如果不够，哪个约束是瓶颈、改进它能带来多大收益？
\end{enumerate}

两层解耦的关键在于：\textbf{一旦外层固定了物理参数，内层剩下的所有耦合都发生在一个共同的物理量上——链路负载。}
本章的核心任务就是证明：无论哪类约束、无论建模精度，内层的可行性条件都可以写成关于这一负载向量的线性不等式。

\begin{center}
\begin{tabular}{l|l}
\hline
\textbf{外层（技术选型）} & \textbf{内层（多约束耦合判断）} \\
\hline
拓扑：Dragonfly $(a,p,h)$, Mesh, Torus, \dots & \\
互连标准：UCIe等级, SerDes速率 & 变量：链路负载向量 \\
裸片：$A$, $P$, $p_{\text{bump}}$ & \\
冷却：风冷/液冷/浸没/微流道 & 性能 + 几何 + 功耗 \\
裸片布局：（由外层决定，非内层优化目标）& 三族线性不等式联立求解 \\
\hline
\end{tabular}
\end{center}
